%%%%%%%%%%%%%%%%%%%%%%%%%%%%%%%%%%%%%%%%%%%%%%%%%%%%%%%%%%%%%%%%%
% PHIẾU HỌC TẬP - SMARTLOGI: TỐI ƯU HÓA VẬN TẢI
% File riêng để in phát cho học sinh
%%%%%%%%%%%%%%%%%%%%%%%%%%%%%%%%%%%%%%%%%%%%%%%%%%%%%%%%%%%%%%%%%
\documentclass[12pt,a4paper]{article}

% --- GÓI CƠ BẢN ---
\usepackage[utf8]{inputenc}
\usepackage[T1]{fontenc}
\usepackage[vietnamese]{babel}
\usepackage{amsmath,amssymb}
\usepackage{array}
\usepackage{booktabs}
\usepackage{tabularx}
\usepackage{graphicx}

% --- FONT ---
\usepackage{newtxtext, newtxmath}
\usepackage[scaled=0.92]{helvet}
\usepackage{microtype}

% --- MÀU SẮC VÀ BỐ CỤC ---
\usepackage[svgnames]{xcolor}
\definecolor{primary}{RGB}{0,102,204}
\definecolor{accent}{RGB}{255,69,0}
\definecolor{success}{RGB}{34,139,34}
\definecolor{TableHeader}{gray}{0.9}

\usepackage[left=1.5cm,right=1.5cm,top=2cm,bottom=2cm,headheight=15pt]{geometry}
\usepackage{fancyhdr}
\pagestyle{fancy}
\fancyhf{}
\renewcommand{\headrule}{\color{primary}\hrule width\textwidth height 1pt}
\fancyhead[L]{\color{primary}\footnotesize\sffamily\textbf{PHIẾU HỌC TẬP - CHỦ ĐỀ STEM: SMARTLOGI}}
\fancyhead[R]{\color{primary}\footnotesize\sffamily\textbf{LỚP 10}}
\fancyfoot[C]{\sffamily\thepage}

% --- GÓI HỖ TRỢ ---
\usepackage{enumitem}
\usepackage{tcolorbox}
\tcbuselibrary{skins, breakable}
\usepackage[colorlinks=true,linkcolor=primary,urlcolor=accent]{hyperref}

% --- MÔI TRƯỜNG ---
\newtcolorbox{phieuhoctap}[1]{
    enhanced, breakable,
    colback=white, colframe=success,
    boxrule=1.5pt, arc=3mm,
    fonttitle=\bfseries\sffamily\Large,
    coltitle=white, colbacktitle=success,
    title=#1,
    before upper={\setlength{\parskip}{0.5em}}
}

\newtcolorbox{tinhuong}{
    enhanced, breakable,
    colback=primary!5, colframe=primary!80,
    boxrule=1pt, arc=2mm,
    fonttitle=\bfseries\sffamily,
    title=\color{accent}TÌNH HUỐNG THỰC TIỄN
}

\newtcolorbox{congthuc}{
    enhanced, breakable,
    colback=accent!5, colframe=accent!80,
    boxrule=1pt, arc=2mm,
    fonttitle=\bfseries\sffamily,
    title=\color{primary}CÔNG THỨC TRỌNG TÂM
}

\newtcolorbox{khaosat}{
    enhanced, breakable,
    colback=gray!5, colframe=gray!60,
    boxrule=1pt, arc=2mm,
    fonttitle=\bfseries\sffamily,
    title=\color{primary}PHIẾU KHẢO SÁT
}

\newcolumntype{C}{>{\centering\arraybackslash}X}
\newcolumntype{L}{>{\raggedright\arraybackslash}X}
\newcommand{\header}[1]{\cellcolor{TableHeader}\sffamily\bfseries#1}

\setlist[itemize,1]{label=\textbullet, leftmargin=*, nosep}
\setlist[itemize,2]{label=--, leftmargin=*, nosep}

\begin{document}

%==============================================================
% PHIẾU HỌC TẬP SỐ 1
%==============================================================
\begin{phieuhoctap}{PHIẾU HỌC TẬP SỐ 1: MÔ HÌNH HÓA BÀI TOÁN LOGISTICS}

\noindent\textbf{Họ và tên:} \underline{\hspace{6cm}} \hfill \textbf{Nhóm:} \underline{\hspace{2cm}}

\vspace{0.5cm}
\begin{tinhuong}
\textbf{Bạn là CEO của startup SmartLogi} - một công ty logistics mới thành lập. Công ty nhận được đơn hàng đầu tiên:
\begin{itemize}
\item Vận chuyển tối thiểu \textbf{9 tấn} hàng hóa
\item Tổng thể tích tối thiểu \textbf{36 m$^3$}
\end{itemize}

\textbf{Bạn có thể thuê 2 loại xe:}
\begin{center}
\begin{tabular}{|l|c|c|c|}
\hline
\textbf{Loại xe} & \textbf{Tải trọng} & \textbf{Thể tích} & \textbf{Giá thuê} \\
\hline
Xe loại A & 3 tấn & 3 m$^3$ & 4 triệu/xe \\
Xe loại B & 1 tấn & 6 m$^3$ & 3 triệu/xe \\
\hline
\end{tabular}
\end{center}

\textbf{Thách thức:} Làm thế nào để chọn phương án vận chuyển với chi phí thấp nhất?
\end{tinhuong}

\vspace{0.5cm}
\textbf{\color{accent}BÀI 1: Xác định ẩn số}
\begin{itemize}
\item Gọi $x$ = số xe loại A cần thuê
\item Gọi $y$ = số xe loại B cần thuê
\item Điều kiện: $x, y \ge 0$ và $x, y \in \mathbb{N}$ (số nguyên không âm)
\end{itemize}

\vspace{0.5cm}
\textbf{\color{accent}BÀI 2: Viết các ràng buộc (điền vào chỗ trống)}

\begin{enumerate}[label=\alph*)]
\item \textbf{Ràng buộc tải trọng:} (Xe A chở 3 tấn, Xe B chở 1 tấn, cần $\ge$ 9 tấn)

\hspace{2cm} $3x + \underline{\hspace{1.5cm}}y \ge \underline{\hspace{1.5cm}}$

\item \textbf{Ràng buộc thể tích:} (Xe A chở 3 m$^3$, Xe B chở 6 m$^3$, cần $\ge$ 36 m$^3$)

\hspace{2cm} $\underline{\hspace{1.5cm}}x + 6y \ge \underline{\hspace{1.5cm}}$

\item \textbf{Giới hạn số xe:}

\hspace{2cm} $0 \le x \le \underline{\hspace{1.5cm}}$ \quad (giả sử tối đa 10 xe A)

\hspace{2cm} $0 \le y \le \underline{\hspace{1.5cm}}$ \quad (giả sử tối đa 9 xe B)
\end{enumerate}

\vspace{0.5cm}
\textbf{\color{accent}BÀI 3: Xác định hàm mục tiêu}
\begin{itemize}
\item Chi phí thuê xe A: \underline{\hspace{2cm}} triệu/xe
\item Chi phí thuê xe B: \underline{\hspace{2cm}} triệu/xe
\item \textbf{Hàm chi phí:} $C = \underline{\hspace{1.5cm}}x + \underline{\hspace{1.5cm}}y \to \min$
\end{itemize}

\vspace{0.5cm}
\textbf{\color{accent}BÀI 4: Dự đoán theo trực giác}

\begin{tabularx}{\textwidth}{|L|C|C|}
\hline
\header{Phương án dự đoán} & \header{Số xe A} & \header{Số xe B} \\
\hline
Phương án của em & \underline{\hspace{2cm}} & \underline{\hspace{2cm}} \\
\hline
\end{tabularx}

\vspace{0.3cm}
\textbf{Kiểm tra tính khả thi:}
\begin{itemize}
\item Chi phí dự đoán: \underline{\hspace{4cm}} triệu VNĐ
\item Tải trọng = \underline{\hspace{2cm}} tấn $\ge$ 9 tấn? \quad $\square$ Đạt \quad $\square$ Không đạt
\item Thể tích = \underline{\hspace{2cm}} m$^3$ $\ge$ 36 m$^3$? \quad $\square$ Đạt \quad $\square$ Không đạt
\end{itemize}

\begin{congthuc}
\textbf{Hệ bất phương trình bậc nhất hai ẩn:}
$$\begin{cases}
3x + y \ge 9 \\
3x + 6y \ge 36 \\
0 \le x \le 10 \\
0 \le y \le 9
\end{cases}$$
\textbf{Hàm mục tiêu:} $C = 4x + 3y \to \min$
\end{congthuc}
\end{phieuhoctap}

\newpage
%==============================================================
% PHIẾU HỌC TẬP SỐ 2
%==============================================================
\begin{phieuhoctap}{PHIẾU HỌC TẬP SỐ 2: HƯỚNG DẪN SỬ DỤNG SMARTLOGI}

\noindent\textbf{Họ và tên:} \underline{\hspace{6cm}} \hfill \textbf{Nhóm:} \underline{\hspace{2cm}}

\vspace{0.5cm}
\textbf{\color{primary}I. GIỚI THIỆU ỨNG DỤNG}
\begin{itemize}
\item \textbf{Tên ứng dụng:} SmartLogi - Tối ưu hóa Vận tải
\item \textbf{Địa chỉ:} File \texttt{SmartLogi.html} (mở bằng trình duyệt) hoặc \href{https://stem-1h8.pages.dev/SmartLogi.html}{stem-1h8.pages.dev/SmartLogi.html}
\item \textbf{Cấu trúc 3 Tab:}
    \begin{itemize}
    \item \textbf{Tab 1 - TÍNH TOÁN:} Nhập tham số, xem miền nghiệm, tìm điểm tối ưu
    \item \textbf{Tab 2 - PHÂN TÍCH KỊCH BẢN:} So sánh kịch bản, phân tích độ nhạy
    \item \textbf{Tab 3 - HỌC TẬP:} Công thức Toán, kiến thức Logistics
    \end{itemize}
\end{itemize}

\vspace{0.5cm}
\textbf{\color{primary}II. NHIỆM VỤ THỰC HÀNH}

\textbf{\color{accent}Bước 1: Mở Tab ``TÍNH TOÁN''}
\begin{itemize}
\item Nhập các tham số theo bài toán gốc:
    \begin{itemize}
    \item \textit{Xe A:} Tải trọng = 3 tấn, Thể tích = 3 m$^3$, Giá = 4 triệu
    \item \textit{Xe B:} Tải trọng = 1 tấn, Thể tích = 6 m$^3$, Giá = 3 triệu
    \item \textit{Yêu cầu:} Tổng tấn = 9, Tổng khối = 36
    \item \textit{Giới hạn:} Max xe A = 10, Max xe B = 9
    \end{itemize}
\end{itemize}

\textbf{Ghi lại kết quả tối ưu:}

\begin{tabularx}{\textwidth}{|L|C|C|C|}
\hline
& \header{Số xe A} & \header{Số xe B} & \header{Chi phí tối ưu} \\
\hline
Kết quả từ ứng dụng & \underline{\hspace{2cm}} & \underline{\hspace{2cm}} & \underline{\hspace{3cm}} triệu \\
\hline
\end{tabularx}

\vspace{0.5cm}
\textbf{\color{accent}Bước 2: Mở Tab ``PHÂN TÍCH KỊCH BẢN''}
\begin{itemize}
\item Nhấn ``+ Thêm kịch bản'' để tạo \textbf{Kịch bản 2} của nhóm em
\item \textit{Thay đổi đề xuất:} (VD: Giá xe A tăng lên 5 triệu)

\underline{\hspace{\textwidth}}
\end{itemize}

\textbf{Ghi lại kết quả mới:}

\begin{tabularx}{\textwidth}{|L|C|C|C|}
\hline
& \header{Số xe A} & \header{Số xe B} & \header{Chi phí} \\
\hline
Kịch bản 2 & \underline{\hspace{2cm}} & \underline{\hspace{2cm}} & \underline{\hspace{3cm}} triệu \\
\hline
\end{tabularx}

\vspace{0.3cm}
\textbf{So sánh:} Quan sát biểu đồ cột. Kịch bản nào tốt hơn? \underline{\hspace{6cm}}

\vspace{0.5cm}
\textbf{\color{accent}Bước 3: Phân tích Độ nhạy (vẫn ở Tab 2)}
\begin{itemize}
\item Kéo xuống xem biểu đồ đường ``Phân tích độ nhạy''
\item \textbf{Nhận xét:} Yếu tố nào làm chi phí thay đổi nhiều nhất?

\underline{\hspace{\textwidth}}

\underline{\hspace{\textwidth}}
\end{itemize}

\vspace{0.5cm}
\textbf{\color{accent}Bước 4: Mở Tab ``HỌC TẬP''}
\begin{itemize}
\item Đọc phần ``Kiến thức Logistics''. Kể tên 2 thuật ngữ mới:
    \begin{enumerate}
    \item \underline{\hspace{6cm}}
    \item \underline{\hspace{6cm}}
    \end{enumerate}
\end{itemize}

\vspace{0.5cm}
\textbf{\color{primary}III. KẾT LUẬN CỦA NHÓM}

\textbf{Phương án tối ưu được nhóm chọn:}

\underline{\hspace{\textwidth}}

\textbf{Lý do lựa chọn:}

\underline{\hspace{\textwidth}}

\underline{\hspace{\textwidth}}
\end{phieuhoctap}

\newpage
%==============================================================
% PHIẾU KHẢO SÁT TRƯỚC KHI HỌC
%==============================================================
\begin{khaosat}

\begin{center}
\Large\bfseries\sffamily PHIẾU KHẢO SÁT TRƯỚC KHI HỌC
\end{center}

\noindent\textbf{Họ và tên:} \underline{\hspace{6cm}} \hfill \textbf{Lớp:} \underline{\hspace{2cm}}

\vspace{0.5cm}
\textbf{Hãy đánh dấu ($\checkmark$) vào ô phù hợp với em:}

\vspace{0.3cm}
\begin{tabularx}{\textwidth}{|L|c|c|c|c|}
\hline
\header{Câu hỏi} & \header{Rất đồng ý} & \header{Đồng ý} & \header{Phân vân} & \header{Không} \\
\hline
1. Em biết Logistics (vận tải) là gì & $\square$ & $\square$ & $\square$ & $\square$ \\
\hline
2. Em đã nghe về ``tối ưu hóa chi phí'' & $\square$ & $\square$ & $\square$ & $\square$ \\
\hline
3. Em biết cách giải Hệ bất phương trình 2 ẩn & $\square$ & $\square$ & $\square$ & $\square$ \\
\hline
4. Em biết cách vẽ miền nghiệm trên mặt phẳng tọa độ & $\square$ & $\square$ & $\square$ & $\square$ \\
\hline
5. Em thấy Toán học có thể ứng dụng vào kinh doanh & $\square$ & $\square$ & $\square$ & $\square$ \\
\hline
6. Em quan tâm đến khởi nghiệp và kinh doanh & $\square$ & $\square$ & $\square$ & $\square$ \\
\hline
\end{tabularx}

\vspace{0.5cm}
\textbf{7. Em mong muốn học được gì từ bài học này?}

\underline{\hspace{\textwidth}}

\underline{\hspace{\textwidth}}
\end{khaosat}

\vspace{1cm}
%==============================================================
% PHIẾU KHẢO SÁT SAU KHI HỌC
%==============================================================
\begin{khaosat}

\begin{center}
\Large\bfseries\sffamily PHIẾU KHẢO SÁT SAU KHI HỌC
\end{center}

\noindent\textbf{Họ và tên:} \underline{\hspace{6cm}} \hfill \textbf{Lớp:} \underline{\hspace{2cm}}

\vspace{0.5cm}
\textbf{Hãy đánh dấu ($\checkmark$) vào ô phù hợp với em:}

\vspace{0.3cm}
\begin{tabularx}{\textwidth}{|L|c|c|c|c|}
\hline
\header{Câu hỏi} & \header{Rất đồng ý} & \header{Đồng ý} & \header{Phân vân} & \header{Không} \\
\hline
1. Em hiểu rõ hơn về bài toán tối ưu hóa & $\square$ & $\square$ & $\square$ & $\square$ \\
\hline
2. Em biết cách mô hình hóa bài toán thực tế thành Hệ BPT & $\square$ & $\square$ & $\square$ & $\square$ \\
\hline
3. Ứng dụng SmartLogi giúp em hiểu bài tốt hơn & $\square$ & $\square$ & $\square$ & $\square$ \\
\hline
4. Em thấy Toán học hữu ích trong kinh doanh thực tế & $\square$ & $\square$ & $\square$ & $\square$ \\
\hline
5. Em muốn học thêm các chủ đề STEM tương tự & $\square$ & $\square$ & $\square$ & $\square$ \\
\hline
6. Em tự tin có thể áp dụng kiến thức này vào thực tế & $\square$ & $\square$ & $\square$ & $\square$ \\
\hline
\end{tabularx}

\vspace{0.5cm}
\textbf{7. Điều em thích nhất về bài học này là gì?}

\underline{\hspace{\textwidth}}

\underline{\hspace{\textwidth}}

\vspace{0.3cm}
\textbf{8. Điều em muốn cải thiện về bài học này là gì?}

\underline{\hspace{\textwidth}}

\underline{\hspace{\textwidth}}
\end{khaosat}

\newpage
%==============================================================
% RUBRIC ĐÁNH GIÁ SẢN PHẨM
%==============================================================
\begin{center}
{\Large\bfseries\sffamily\color{primary} RUBRIC ĐÁNH GIÁ SẢN PHẨM}
\end{center}

\vspace{0.3cm}
\noindent\textbf{Tên nhóm:} \underline{\hspace{4cm}} \hfill \textbf{Người đánh giá:} \underline{\hspace{4cm}}

\vspace{0.5cm}
\begin{tabularx}{\textwidth}{|>{\raggedright\arraybackslash}p{2.8cm}|>{\raggedright\arraybackslash}X|>{\raggedright\arraybackslash}X|>{\raggedright\arraybackslash}X|c|}
\hline
\rowcolor{TableHeader}
\sffamily\bfseries Tiêu chí & \sffamily\bfseries Xuất sắc (4đ) & \sffamily\bfseries Khá (3đ) & \sffamily\bfseries Cần cải thiện (1-2đ) & \sffamily\bfseries Điểm \\
\hline

\textbf{1. Mô hình hóa Toán học} 
& 
\begin{itemize}[leftmargin=*,noitemsep,topsep=0pt]
    \item Hệ BPT chính xác
    \item Xác định đúng hàm mục tiêu
    \item Giải thích rõ ràng
\end{itemize}
& 
\begin{itemize}[leftmargin=*,noitemsep,topsep=0pt]
    \item Hệ BPT cơ bản đúng
    \item Có sai sót nhỏ
\end{itemize}
& 
\begin{itemize}[leftmargin=*,noitemsep,topsep=0pt]
    \item Hệ BPT sai
    \item Không xác định được mục tiêu
\end{itemize}
& \underline{\hspace{1cm}} \\
\hline

\textbf{2. Sử dụng công nghệ}
& 
\begin{itemize}[leftmargin=*,noitemsep,topsep=0pt]
    \item Thành thạo cả 3 tabs
    \item Phân tích độ nhạy đúng
\end{itemize}
& 
\begin{itemize}[leftmargin=*,noitemsep,topsep=0pt]
    \item Sử dụng tốt Tab 1, 2
    \item Còn lúng túng Tab 3
\end{itemize}
& 
\begin{itemize}[leftmargin=*,noitemsep,topsep=0pt]
    \item Chỉ dùng được Tab 1
    \item Không phân tích được
\end{itemize}
& \underline{\hspace{1cm}} \\
\hline

\textbf{3. Tính khả thi phương án}
& 
\begin{itemize}[leftmargin=*,noitemsep,topsep=0pt]
    \item Đề xuất 2+ kịch bản khả thi
    \item Phân tích ưu nhược chi tiết
\end{itemize}
& 
\begin{itemize}[leftmargin=*,noitemsep,topsep=0pt]
    \item Đề xuất 1-2 kịch bản
    \item Phân tích cơ bản
\end{itemize}
& 
\begin{itemize}[leftmargin=*,noitemsep,topsep=0pt]
    \item Phương án không khả thi
    \item Thiếu phân tích
\end{itemize}
& \underline{\hspace{1cm}} \\
\hline

\textbf{4. Trình bày và bảo vệ}
& 
\begin{itemize}[leftmargin=*,noitemsep,topsep=0pt]
    \item Mạch lạc, logic
    \item Trực quan hiệu quả
    \item Trả lời trọn vẹn
\end{itemize}
& 
\begin{itemize}[leftmargin=*,noitemsep,topsep=0pt]
    \item Trình bày rõ ràng
    \item Có hỗ trợ trực quan
\end{itemize}
& 
\begin{itemize}[leftmargin=*,noitemsep,topsep=0pt]
    \item Trình bày rời rạc
    \item Không trả lời được
\end{itemize}
& \underline{\hspace{1cm}} \\
\hline
\end{tabularx}

\vspace{0.5cm}
\noindent\textbf{TỔNG ĐIỂM:} \underline{\hspace{2cm}} / 16 điểm

\vspace{0.5cm}
\noindent\textbf{Nhận xét của giáo viên:}

\underline{\hspace{\textwidth}}

\underline{\hspace{\textwidth}}

\underline{\hspace{\textwidth}}

\vspace{1cm}
%==============================================================
% PHIẾU ĐÁNH GIÁ ĐỒNG CẤP
%==============================================================
\begin{center}
{\Large\bfseries\sffamily\color{primary} PHIẾU ĐÁNH GIÁ ĐỒNG CẤP (PEER ASSESSMENT)}
\end{center}

\vspace{0.3cm}
\noindent\textbf{Tên em:} \underline{\hspace{4cm}} \hfill \textbf{Nhóm:} \underline{\hspace{2cm}}

\vspace{0.5cm}
\textbf{Đánh giá các thành viên trong nhóm} (thang điểm 1-5: 1=Yếu, 5=Xuất sắc)

\vspace{0.3cm}
\begin{tabularx}{\textwidth}{|L|c|c|c|c|}
\hline
\header{Tên thành viên} & \header{Ý tưởng} & \header{Hợp tác} & \header{Hoàn thành} & \header{Tổng} \\
\hline
\underline{\hspace{4cm}} & \underline{\hspace{1cm}} & \underline{\hspace{1cm}} & \underline{\hspace{1cm}} & \underline{\hspace{1cm}} \\
\hline
\underline{\hspace{4cm}} & \underline{\hspace{1cm}} & \underline{\hspace{1cm}} & \underline{\hspace{1cm}} & \underline{\hspace{1cm}} \\
\hline
\underline{\hspace{4cm}} & \underline{\hspace{1cm}} & \underline{\hspace{1cm}} & \underline{\hspace{1cm}} & \underline{\hspace{1cm}} \\
\hline
\underline{\hspace{4cm}} & \underline{\hspace{1cm}} & \underline{\hspace{1cm}} & \underline{\hspace{1cm}} & \underline{\hspace{1cm}} \\
\hline
\end{tabularx}

\vspace{0.5cm}
\noindent\textbf{Thành viên đóng góp nhiều nhất:} \underline{\hspace{6cm}}

\noindent\textbf{Lý do:} \underline{\hspace{\textwidth}}

\end{document}
